\documentclass[11pt]{article}

\usepackage[margin=1in]{geometry}
\usepackage{amsmath,amssymb,amsthm}
\usepackage{enumitem}
\usepackage{hyperref}
\usepackage{microtype}

\hypersetup{
  colorlinks=true,
  linkcolor=black,
  urlcolor=blue,
  citecolor=black
}

\title{\textbf{Research Warm-Up Project}\\[0.4em]
\large Adaptive Coverage Policies in Conformal Prediction}
\author{}
\date{}

\begin{document}

\maketitle
\thispagestyle{empty}

\bigskip

\section{Overview}

The goal of this project is to provide hands-on exposure to modern conformal
prediction methods and to develop the skills needed to read, reproduce, and
extend recent machine learning research.

The project is centered around the paper

\begin{quote}
  Gauthier, Bach, and Jordan,
  \emph{Adaptive Coverage Policies in Conformal Prediction},
  AISTATS 2026. \texttt{arXiv:2510.04318}
\end{quote}

\noindent By the end of this project, you should be able to:

\begin{itemize}
  \item Implement split conformal prediction from scratch.
  \item Understand prediction sets and marginal coverage guarantees.
  \item Understand the role of e-values in enabling post-hoc valid inference.
  \item Reproduce the methodology proposed in the ACP paper.
  \item Reproduce at least one experimental result from the paper.
  \item Conduct a small extension study.
  \item Communicate results in the form of a short research report.
\end{itemize}

Throughout the project, maintain a Git repository containing all code,
experiments, and figures. Every figure appearing in your report should be
reproducible from the repository.

% -----------------------------------------------------------------------
\section{Part 1: Classical Conformal Prediction}
% -----------------------------------------------------------------------

\subsection{Dataset and Model}

Use the CIFAR-10 dataset, which is the dataset used in the ACP paper for
classification. Train an image classifier on the full CIFAR-10 training set
and treat it as a black-box predictor throughout the project. Use the same
train/calibration/test split throughout.

\subsection{Split Conformal Prediction}

Implement split conformal prediction from scratch.

Let $\hat p_y(x)$ denote the predicted probability (softmax output) assigned
to class $y$. Use the softmax conformity score
\[
  S(x,y) \;=\; 1 - \hat p_y(x).
\]
Note that this score is non-negative and \emph{negatively oriented}: lower
scores indicate better predictions. This property will matter when we later
switch to e-value-based prediction sets, which require exactly these two
conditions.

Construct conformal prediction sets at a fixed miscoverage level $\alpha$
using the standard empirical-quantile threshold: include all labels $y$
whose score does not exceed the $\lceil(1-\alpha)(n+1)\rceil/n$ empirical
quantile of the $n$ calibration scores.

You may use existing libraries for training the classifier, but the conformal
prediction procedure itself should be implemented by you.

\subsection{Experiments}

Evaluate:
\begin{itemize}
  \item Empirical coverage.
  \item Average prediction set size.
  \item Distribution of prediction set sizes.
\end{itemize}
Repeat for several values of $\alpha$.

\subsection{Questions}

\begin{enumerate}[label=\arabic*.]
  \item How does prediction set size vary with $\alpha$?
  \item Which examples tend to receive the largest prediction sets?
  \item Which examples tend to receive the smallest prediction sets?
  \item What practical limitations does a fixed, global $\alpha$ create?
\end{enumerate}

% -----------------------------------------------------------------------
\section{Part 2: Understanding the ACP Paper}
% -----------------------------------------------------------------------

Read the paper carefully before implementing any part of the method. Your
goal is not merely to reproduce code, but to understand the main ideas.
Prepare written answers to the following questions.

\subsection{Reading Questions}

\begin{enumerate}[label=\arabic*.]

\item What problem is ACP attempting to solve, and why does a globally fixed
$\alpha$ lead to suboptimal or uninformative predictions?

\item What is an e-variable? Write out the definition. What is the
\emph{soft-rank e-variable} used in the paper, and how does it differ from
the rank-based comparison used in classical conformal prediction?

\item How is the conformal e-prediction set (eq.~4 in the paper) constructed
from the soft-rank e-variable? How does this differ from the
quantile-threshold construction you implemented in Part~1?

\item State the post-hoc validity guarantee (Proposition~2.2) precisely.
How does it differ from the standard marginal coverage guarantee
$\mathbb{P}(Y_{\mathrm{test}} \in \hat C_n^\alpha(X_{\mathrm{test}})) \ge
1-\alpha$?

\item Why does post-hoc validity allow the miscoverage level $\tilde\alpha$
to be chosen in a data-dependent manner, whereas classical conformal
prediction does not permit this?

\item What is a coverage policy, and how does the paper formally define it?
What inputs does it take, and why does the formal definition reduce to the
\emph{sum} of calibration scores rather than the full vector?

\item The coverage policy is trained using a leave-one-out (LOO) procedure
on the calibration set. What is the purpose of the LOO procedure here?
(\emph{Hint}: it is not primarily about avoiding overfitting or reusing data
efficiently, but about constructing training data.)

\item What is the training objective $\mathcal{L}_\lambda(\theta)$ (eq.~9)?
Describe each of its two terms and what each encourages the policy to do.
What is the role of the regularization parameter $\lambda$, and how does the
paper recommend choosing it?

\item The paper uses the \emph{cross-entropy} as the conformity score in its
classification experiments:
$S(x, y) = -\log \hat p_y(x)$.
What two properties of a score function are required for the soft-rank
e-variable to be a valid e-value? Verify that the cross-entropy score
satisfies them.

\item What assumptions are required for the post-hoc validity guarantee?

\end{enumerate}

% -----------------------------------------------------------------------
\section{Part 3: Implementing Adaptive Coverage Policies}
% -----------------------------------------------------------------------

Implement the ACP methodology described in the paper. The implementation
should include:

\begin{itemize}

  \item The \textbf{conformity score}. Use the cross-entropy score
  $S(x,y) = -\log \hat p_y(x)$ used in the paper (not the softmax score
  from Part~1).

  \item The \textbf{soft-rank e-variable} and the corresponding
  \textbf{conformal e-prediction set} (eqs.~2 and~4 in the paper). This is
  the core prediction-set construction; it replaces the quantile-based
  construction used in Part~1.

  \item The \textbf{leave-one-out procedure} for generating pseudo
  calibration--test pairs from the calibration set.

  \item A \textbf{neural network coverage policy} $\tilde\alpha_\theta$
  trained by minimizing the objective $\mathcal{L}_\lambda(\theta)$, using
  the smooth sigmoid surrogate for the non-differentiable set-size term
  (eq.~5 in the paper).

  \item A \textbf{procedure for selecting $\lambda$} to target a desired
  expected prediction set size at test time (Algorithm~2).

  \item \textbf{Test-time inference}: computing $\tilde\alpha_\theta^{\mathrm{test}}$
  and constructing the corresponding prediction set.

\end{itemize}

Do not consult the authors' implementation until after you have completed
your own. Document any decisions not fully specified in the paper (e.g.\
network architecture, sharpness parameter $k$, optimizer and learning rate).

\subsection{Verification}

Before running full experiments, verify that:
\begin{itemize}
  \item The trained policy produces non-trivial variation in miscoverage
  levels $\tilde\alpha$ across test examples.
  \item The distribution of $\tilde\alpha$ values shifts as $\lambda$
  changes.
  \item Empirical coverage is consistent with the post-hoc validity
  guarantee (Proposition~2.2).
\end{itemize}

% -----------------------------------------------------------------------
\section{Part 4: Reproducing a Published Result}
% -----------------------------------------------------------------------

Select one figure or table from the paper's experiments and reproduce it
as faithfully as possible. Good candidates include:

\begin{itemize}
  \item Training dynamics of the coverage policy under different values
  of $\lambda$ (Figure~1 in the paper).
  \item Average prediction set size as a function of $\lambda$.
  \item Empirical miscoverage as a function of $\lambda$.
  \item The distribution of learned miscoverage levels $\tilde\alpha$
  across test examples under different $\lambda$.
\end{itemize}

\subsection{Requirements}

Use:
\begin{itemize}
  \item The same dataset (CIFAR-10).
  \item The same score function (cross-entropy).
  \item The same train/calibration/test protocol described in the paper.
  \item The same evaluation metrics.
\end{itemize}

\subsection{Deliverables}

\begin{enumerate}[label=\arabic*.]
  \item The reproduced figure.
  \item A one-page implementation summary.
  \item A discussion of any discrepancies.
\end{enumerate}

% -----------------------------------------------------------------------
\section{Part 5: Extension Study}
% -----------------------------------------------------------------------

The ACP framework trains a neural coverage policy on top of a
conformal e-predictor. An important open question is whether its benefits
depend on the choice of conformity score.

In this part, investigate the interaction between ACP and two different
conformity scores. Recall that the soft-rank e-variable requires scores to
be \emph{non-negative and negatively oriented} (lower scores must indicate
better predictions). Before plugging any score into the e-value construction,
verify these two properties explicitly.

\subsection{Candidate Scores}

\begin{enumerate}

\item \textbf{Cross-entropy score} (used in the paper):
\[
  S(x,y) = -\log \hat p_y(x).
\]

\item \textbf{APS score} (Romano et al., 2020): Sort classes in decreasing
order of predicted probability; $S(x,y)$ is the cumulative predicted
probability up to and including class $y$. Verify non-negativity and
negative orientation before use.

\end{enumerate}

\subsection{Experiments}

For each conformity score:
\begin{itemize}
  \item Train an ACP policy across several values of $\lambda$.
  \item Measure average prediction set size and compare to the fixed-$\alpha$
  conformal e-predictor baseline (without ACP).
  \item Verify empirical coverage against the post-hoc validity guarantee.
  \item Plot the distribution of learned miscoverage levels $\tilde\alpha$
  across test examples.
\end{itemize}

\subsection{Research Questions}

\begin{enumerate}[label=\arabic*.]
  \item Does ACP consistently reduce average prediction set size regardless
  of the conformity score?
  \item Does the trained policy produce a similar distribution of
  $\tilde\alpha$ values for both scores?
  \item Which conformity score yields the largest efficiency gains from ACP?
  \item Are the paper's conclusions robust to the choice of conformity score?
  \item What properties of a conformity score appear to make it well-suited
  for use within the ACP framework?
\end{enumerate}

% -----------------------------------------------------------------------
\section{Final Report}
% -----------------------------------------------------------------------

Write a short report of approximately 3--5 pages, containing:

\begin{enumerate}[label=\arabic*.]
  \item Introduction.
  \item Background on conformal prediction, e-values, and post-hoc validity.
  \item Summary of the ACP methodology, including the soft-rank e-variable,
  the training objective, and the leave-one-out procedure.
  \item Reproduction study.
  \item Extension study.
  \item Conclusions.
\end{enumerate}

Focus on clear explanations, careful experimentation, and thoughtful
interpretation of results.

% -----------------------------------------------------------------------
\section{Stretch Goal}
% -----------------------------------------------------------------------

For students interested in theory, study the validity guarantees of both
split conformal prediction and ACP.

\begin{enumerate}

\item Explain why split conformal prediction achieves finite-sample marginal
coverage under exchangeability.

\item The soft-rank e-variable is:
\[
  E = \frac{S(X_{\mathrm{test}},Y_{\mathrm{test}})}
           {\frac{1}{n+1}\!\left(\sum_{i=1}^{n} S(X_i,Y_i)
            + S(X_{\mathrm{test}},Y_{\mathrm{test}})\right)}.
\]
Verify that $\mathbb{E}[E] \le 1$ when scores are exchangeable and
non-negative, thereby confirming it is a valid e-variable.
(\emph{Hint}: use exchangeability to argue that the test score is a uniform
draw from the $n+1$ scores, then compute the expectation directly.)

\item State the post-hoc validity guarantee (Proposition~2.2) precisely.
Why does this guarantee hold even when $\tilde\alpha$ depends on all the
data, whereas the standard guarantee (eq.~1) requires $\alpha$ to be fixed
in advance?

\item The ACP guarantee is
\[
  \mathbb{E}\!\left[\frac{\mathbb{P}(Y_{\mathrm{test}} \notin
  \hat C_n^{\tilde\alpha}(X_{\mathrm{test}}) \mid \tilde\alpha)}
  {\tilde\alpha}\right] \le 1,
\]
whereas the standard guarantee is
$\mathbb{P}(Y_{\mathrm{test}} \notin \hat C_n^\alpha) \le \alpha$.
Explain the difference between these two statements. Under what condition
do they coincide?

\end{enumerate}

A complete proof is not required. The goal is to understand the structure
of the arguments and the role of e-values in enabling post-hoc valid
adaptive inference.

\end{document}